\section{Limitations and questions for the ELVES/JAM authors}
\label{sec:limitations}

The purpose of this paper is to expose assumptions, not to claim a complete security analysis. Several issues require author or implementation review.

\subsection{Correlated honest failure is outside the ELVES proof model}

The branching extension assumes that a buggy honest auditor behaves like an auditor that cannot produce a useful negative judgment for fault \(x\). This is a natural substitution in the reproduction mean, but it is not proved by ELVES. The exact coupling to adaptive crashes, tranche timing, and the paper's stronger committee-size concentration bounds should be checked by the ELVES authors.

\subsection{The verdict analysis assumes enough judgments become available}

The finite-population verdict conditions in \Cref{sec:attribution} concern which signed judgment compositions can be assembled. If too many validators are offline, a verdict may fail to reach \(K\) total judgments at all. That is a distinct liveness failure.

\subsection{What does the staking layer do with wonky?}

This is the highest-priority clarification.

The Gray Paper records wonky reports in a dedicated set and removes non-positive reports from the availability path. Its ordinary culprit/fault rules refer to good or bad reports, and its punish set is populated by those offense extrinsics. However, the chain's validator rewards and higher-level economic policy are delegated.

Therefore the paper's narrow claim is:
\begin{quote}
A wonky verdict does not create the ordinary Gray-Paper culprit/fault offender record.
\end{quote}
It does \emph{not} follow that a production staking system must impose zero cost. If the staking layer explicitly penalizes guarantors or judges associated with wonky, the economic attribution gap narrows.

\subsection{The one-vote finite-population gap}

At \(N=1023\), a global composition of 682 positive and 341 negative judgments does not appear to contain a valid 683-judgment subset with any of the three permitted positive counts \(683\), \(341\), or \(0\). This may be resolved elsewhere in the intended protocol semantics, or it may be an unreachable/irrelevant edge case under assumptions not represented in our abstraction. It should be confirmed rather than silently approximated by the continuous \(2/3\) boundary.

\subsection{Strategic entry is not established}

The probability that at least two of three guarantors share a fault does not imply an attacker can freely choose such a core. Strategic exploitation requires a mapping from attacker-controlled work submission to core assignment and timing. The manuscript therefore separates:
\begin{itemize}
\item \emph{accidental correlated failure}, where assignment probabilities enter the event probability; and
\item \emph{strategic exploitation}, where an attacker conditions on known favorable assignments only if the protocol gives it sufficient choice.
\end{itemize}

\subsection{Client identity is not fault identity}

A client-diversity rule can miss shared dependencies and can be gamed when identity is self-declared. Conversely, two versions of the same client might not share a particular bug. The model's fundamental variable is fault-domain membership, which is generally latent until a fault is discovered.

\subsection{Questions to send with the paper}

We would ask the ELVES and JAM authors five concrete questions:
\begin{enumerate}[label=Q\arabic*.]
\item Is replacing the honest offspring share \((1-\gamma)\) by a fault-specific independently correct share \(h_x\) valid under the ELVES adaptive-crash proof, or does another part of the proof change?
\item Under the current JAM dispute semantics, is our reading correct that a wonky verdict creates no culprit/fault offender records? Does the intended staking layer nevertheless act on the wonky set?
\item Is the \(N=1023\), 682-positive/341-negative verdict composition handled by another rule, or is the literal three-count verdict restriction intentional?
\item When validator count changes while the 341-core constant remains fixed, is the expected initial audit count per report intended to scale as \(10N/341\) once all cores are active?
\item Are implementation/fault-domain diversity constraints considered anywhere in guarantor assignment, validator selection, or staking beyond the ecosystem-level encouragement of multiple clients?
\end{enumerate}

A negative answer to any of these can change the conclusions. That is precisely why the results are presented as a reviewable extension rather than a vulnerability claim.
